\documentclass[12pt]{article}

\usepackage[utf8]{inputenc}
\usepackage[T1]{fontenc}
\usepackage{amsmath, amssymb, amsthm}
\usepackage{geometry}
\geometry{a4paper, margin=2.5cm}

\title{Omotopie și Domenii Simplu Conexe}
\author{}
\date{}

\theoremstyle{definition}
\newtheorem{definition}{Definiție}

\theoremstyle{plain}
\newtheorem{theorem}{Teoremă}

\begin{document}

\maketitle

\begin{definition}
(\textbf{Omotopia a două curbe})  
Fie $\gamma_0, \gamma_1 : [a,b] \to \Omega$ două curbe cu aceleași capete, adică



\[
\gamma_0(a)=\gamma_1(a), \qquad \gamma_0(b)=\gamma_1(b).
\]



Spunem că $\gamma_0$ și $\gamma_1$ sunt omotope în $\Omega$ dacă există o funcție continuă



\[
H : [a,b] \times [0,1] \to \Omega
\]



astfel încât



\[
H(t,0)=\gamma_0(t), \qquad H(t,1)=\gamma_1(t),
\]



iar pentru fiecare $s\in[0,1]$, curba $t \mapsto H(t,s)$ are aceleași capete.
\end{definition}

\begin{definition}
(\textbf{Curbă închisă omotopă cu un punct})  
O curbă închisă $\gamma : [a,b] \to \Omega$ se numește omotopă cu un punct în $\Omega$ dacă există o omotopie



\[
H : [a,b] \times [0,1] \to \Omega
\]



astfel încât



\[
H(t,0)=\gamma(t), \qquad H(t,1)=z_0,
\]



unde $z_0$ este un punct fix din $\Omega$.  
În acest caz, curba poate fi contractată continuu până la punctul $z_0$ fără a părăsi domeniul.
\end{definition}

\begin{definition}
(\textbf{Mulțime simplu conexă})  
Un domeniu $\Omega \subset \mathbb{C}$ se numește simplu conex dacă orice curbă închisă din $\Omega$ este omotopă cu un punct în $\Omega$.  
Echivalent: orice buclă din $\Omega$ poate fi contractată continuu la un punct fără a ieși din domeniu.
\end{definition}

\begin{theorem}
\textbf{Teorema lui Cauchy pe domenii simplu conexe.}  
Fie $\Omega$ un domeniu simplu conex și $f:\Omega \to \mathbb{C}$ o funcție olomorfă.  
Dacă $\gamma_0$ și $\gamma_1$ sunt două curbe omotope în $\Omega$ și au aceleași capete, atunci



\[
\int_{\gamma_0} f(z)\,dz = \int_{\gamma_1} f(z)\,dz.
\]



În special, pentru orice curbă închisă $\gamma$ din $\Omega$ are loc



\[
\int_{\gamma} f(z)\,dz = 0.
\]


\end{theorem}

\end{document}
